\section{Bellman-Ford Algorithm}
\label{sec:graph-bellman-ford}

\problemstatement{Given a weighted directed graph, possibly containing
\textbf{negative} edge weights, find the shortest distance from a
source to every other node -- or detect that a \textbf{negative-weight
cycle} makes the notion of ``shortest path'' undefined.}

\begin{patternbox}
\emph{``Shortest path, negative weights allowed''} is solved by
relaxing \textbf{every edge}, $V{-}1$ times over (no priority queue, no
greedy frontier choice -- Dijkstra's greedy finalization is provably
unsafe the moment a negative edge exists, as discussed in
Section~\ref{sec:graph-dijkstra}). A final, $V$-th round that still
finds an improvement reveals a negative cycle.
\end{patternbox}

\subsection*{Why $V-1$ rounds are exactly enough}

In a graph with no negative cycle, the shortest path between any two
nodes is always a \textbf{simple} path (visits no node twice) -- a path
revisiting a node would include a cycle, and removing that cycle could
only shorten the path further unless the cycle's total weight is
\emph{exactly} $0$, in which case removing it doesn't change the
length either way. A simple path in a graph with $V$ nodes has at most
$V{-}1$ edges. Each full round of relaxing every edge extends the
\emph{confirmed-correct} prefix of every shortest path by (at least)
one more edge -- so after $V{-}1$ rounds, every shortest path (having
at most $V{-}1$ edges) is guaranteed fully relaxed and correct.

\subsection*{Why a $V$-th round detects a negative cycle}

If, after $V{-}1$ rounds (when every legitimate shortest path must
already be finalized), a $V$-th round of relaxation \emph{still} finds
an improvement on some edge, that improvement cannot correspond to any
genuine shortest path (those are already settled) -- it can only be
explained by a cycle whose total weight is negative, allowing
``distance'' to be driven arbitrarily low by looping through it
repeatedly.

\subsection*{Algorithm}

\begin{enumerate}
  \item Initialize $\textit{dist}[\textit{src}]=0$, all others
        $\infty$.
  \item Repeat $V{-}1$ times: for every edge $(u,v,w)$, if
        $\textit{dist}[u] \neq \infty$ and $\textit{dist}[u]+w <
        \textit{dist}[v]$, update $\textit{dist}[v]$.
  \item Run one more relaxation pass over all edges; if any edge can
        still be relaxed, report a negative cycle.
\end{enumerate}

\subsection*{C++ implementation}

\begin{lstlisting}
#include <bits/stdc++.h>
using namespace std;

vector<int> bellmanFord(int n, vector<vector<int>>& edges, int src) {
    vector<int> dist(n, INT_MAX);
    dist[src] = 0;

    for (int i = 0; i < n - 1; i++) {
        for (auto& e : edges) {
            int u = e[0], v = e[1], w = e[2];
            if (dist[u] != INT_MAX && dist[u] + w < dist[v]) {
                dist[v] = dist[u] + w;
            }
        }
    }

    for (auto& e : edges) {
        int u = e[0], v = e[1], w = e[2];
        if (dist[u] != INT_MAX && dist[u] + w < dist[v]) {
            return {-1}; // sentinel: negative cycle detected
        }
    }
    return dist;
}

int main() {
    vector<vector<int>> edges = {{0,1,4},{0,2,5},{1,2,-3},{2,3,4}};
    vector<int> dist = bellmanFord(4, edges, 0);
    for (int d : dist) cout << d << " ";
    cout << endl; // 0 4 1 5
    return 0;
}
\end{lstlisting}

\subsection*{Dry run}

\dryrun{Take edges $0{\to}1(4)$, $0{\to}2(5)$, $1{\to}2(-3)$,
$2{\to}3(4)$, $n=4$, source $0$.}

$\textit{dist}=[0,\infty,\infty,\infty]$. Round $1$: relax $0{\to}1$
($\infty\to4$); relax $0{\to}2$ ($\infty\to5$); relax $1{\to}2$
($4{-}3=1<5$, update to $1$); relax $2{\to}3$ ($1+4=5<\infty$, update
to $5$). $\textit{dist}=[0,4,1,5]$. Round $2$: relax $0{\to}1$ ($4$
not $<4$); relax $0{\to}2$ ($5$ not $<1$); relax $1{\to}2$ ($1$ not
$<1$); relax $2{\to}3$ ($5$ not $<5$) -- no changes,
\textbf{stabilized} early (rounds can stop as soon as a full pass makes
no change, though the algorithm above runs the full $V{-}1$ rounds
regardless). Round $3$ ($n{-}1=3$ total): no changes. Final check
round: no edge relaxes further -- no negative cycle. Final:
$[0,4,1,5]$.

\begin{complexitybox}
\textbf{Time Complexity.} $O(V \cdot E)$ -- $V{-}1$ rounds (plus one
check round), each scanning all $E$ edges.
\textbf{Space Complexity.} $O(V)$ for \textit{dist}.
\end{complexitybox}

\begin{takeawaybox}
Bellman-Ford trades Dijkstra's $O((V+E)\log V)$ speed for two things
Dijkstra cannot offer: correctness with negative edges, and the ability
to detect negative cycles at all. When a problem's weights are
guaranteed non-negative, Dijkstra remains the faster default --
Bellman-Ford is the fallback specifically for when that guarantee is
unavailable.
\end{takeawaybox}
